\paragraph{Shared-query sensitivity to stored-reading errors.}
Across all 279 external blocks, the frozen shared policy guarantees that
all six normalized comparison widths are at most $1/2$ on 165 blocks under
fixed internal-reading errors bounded by $1\%$ of block energy; the best
deterministic batch does so on 120.  On the 45 additional blocks,
feasible fixed-record certificates prove that every randomized batch has
worst-case precision success below one.  The certified batch success
upper bound has minimum and median $1/2$, interquartile range
$[1/2,2/3]$, and maximum $4/5$; 26 bounds are at most $1/2$ and 37 are
at most $0.7$.  These are upper bounds on the unrestricted batch optimum,
not claims that each bound is attained.  The 45 blocks span all three
external systems (11, 11, and 23 blocks).

Capacity information reduces the benefit.  With a per-hour cap of half
the block energy as well as the same $1\%$ error allowance, a policy
frozen from the capacity-aware noiseless planner certifies 213 of the
279 coefficient instances, versus 204 for the best deterministic batch;
nine have a certified randomized-batch limitation, with upper bounds
between $2/3$ and $7/8$.  Only 200 recorded profiles satisfy this illustrative
cap.  Within that compatible subset the counts are 153 versus 147, with
six separations.  The original uncapped policy certifies only 191 of the
279 capped instances, which motivates using declared capacity information
when planning.  Both policies are fixed before the fine readings are
opened; neither is claimed optimal for bounded-error acquisition.

The comparison uses two shared readings, uncapped noiseless initial widths for
normalization, exact block endpoints, and the same permitted queries for
every method.  A separately implemented rational checker verifies the
saved contrast coefficients, branch coverage, every reported precision
status, and the complete fixed-error adverse records.  Supplementary
methods (Online Resource~1, S8) give the certificate construction.
